% Appendix: theoretical foundations (v1, 2026-08-20). Paper form of the hub theory program
% (chunking-theory Part III, three-forces, adaptive-exec-map claim-strength guard).
% Assumed preamble (main.tex): amsmath, booktabs, natbib, amsthm with
%   \newtheorem{theorem}{Theorem}  \newtheorem{proposition}{Proposition}
%   \newtheorem{lemma}{Lemma}      \newtheorem{corollary}{Corollary}
%   \newtheorem{fact}{Fact}        \theoremstyle{remark}\newtheorem{remark}{Remark}
% Style: no em-dashes, no prose colons. Imported results are stated as Facts with citations;
% everything labeled Theorem/Proposition/Lemma/Corollary here is ours.

\section{Commitment as an Optimization Dimension}
\label{app:theory}

This appendix formalizes the claims of Section~\ref{sec:intro}. We show that
(i)~valuing commitments of different lengths requires semi-Markov bookkeeping, without which the
critic acquires a spurious preference for long commitments (Proposition~\ref{prop:bookkeeping});
(ii)~deployment-time selection is bounded by the support of the frozen policy and cannot lengthen
commitments (Proposition~\ref{prop:selection}, Lemma~\ref{lem:fullchunk});
(iii)~the gain available to an adaptive commitment decomposes into a policy-dependent part that
policy improvement absorbs and an irreducible part set by the environment
(Theorems~\ref{thm:decomp}--\ref{thm:floor}); and
(iv)~closed-loop demonstrations induce an optimism that grows with the commitment length and that
correct bookkeeping does not remove (Proposition~\ref{prop:leakage}).
Together these results locate the mechanism our method needs. The policy must be improved in both
of its output dimensions, and the two biases that act on the commitment must be handled by
different tools.

\subsection{Setup}
\label{app:setup}

Let $M=(\mathcal S,\mathcal A,T,r,\gamma)$ be a Markov decision process with rewards in $[0,1]$
and discount $\gamma\in(0,1)$. A \emph{chunk policy} with horizon $H$ is a map
$\pi\colon\mathcal S\to\Delta(\mathcal A^{H})$, and a \emph{commitment selector} is a map
$\kappa\colon\mathcal S\to\{1,\dots,H\}$. The pair $(\pi,\kappa)$ is executed by querying
$a_{1:H}\sim\pi(\cdot\mid s)$, executing the first $\kappa(s)$ actions open loop, and requerying at
the state reached; we write $V^{\pi,\kappa}$ for the value of this process and abbreviate
$V^{\pi,k}$ when $\kappa\equiv k$ is constant.

Requery times are decision points of a semi-Markov decision process whose actions are the
temporally extended pairs $(a_{1:k},k)$~\citep{sutton1999options, bradtke1994smdp}. The option
value of committing to the first $k$ actions of a chunk and thereafter following $(\pi,\kappa)$ is
\begin{equation}
\label{eq:smdp-backup}
Q^{\pi,\kappa}(s, a_{1:H}, k)
  \;=\;
  \mathbb E\!\left[\;\sum_{j=0}^{k-1}\gamma^{j}\, r(s_j,a_{j+1})
  \;+\;\gamma^{k}\, V^{\pi,\kappa}(s_k)\right],
\end{equation}
where the expectation is over the open-loop rollout of $a_{1:k}$ from $s_0=s$. Chunk-level
temporal-difference learning instantiates \eqref{eq:smdp-backup} for fixed $k=H$~\citep{qc}, and
dual-discount schemes are its within-chunk and across-chunk factorization~\citep{deas}. Because the
$H$ prefixes of one chunk share a single rollout, a causal critic over the chunk emits all $H$
option values in one pass; this is the object our method trains.

Three optimal values will be compared,
\[
V^{\star}_{1}\;:=\;\sup_{\text{closed-loop}}V,
\qquad
V^{\star}_{H}\;:=\;\sup_{\pi}V^{\pi,H},
\qquad
V^{\star}_{\mathrm{ada}}\;:=\;\sup_{\pi,\kappa}V^{\pi,\kappa}.
\]

\begin{lemma}[Sandwich]
\label{lem:sandwich}
$V^{\star}_{H}\le V^{\star}_{\mathrm{ada}}\le V^{\star}_{1}$.
\end{lemma}
\begin{proof}
The left inequality holds because $\kappa\equiv H$ is an admissible selector. For the right one,
any execution of $(\pi,\kappa)$ is a particular policy measurable with respect to the observation
filtration, so its value cannot exceed the closed-loop optimum.
\end{proof}

\subsection{Bookkeeping. Comparing commitments requires the semi-Markov backup}
\label{app:bookkeeping}

Consider the \emph{duration-blind} backup that charges one step of discount per decision
regardless of the committed duration,
$\widetilde Q(s,a_{1:H},k):=\mathbb E[\sum_{j<k}\gamma^{j}r_j+\gamma\,V(s_k)]$. This is the most
favorable form of the error, since the within-commitment rewards are still discounted correctly;
only the bootstrap term is mis-charged.

\begin{proposition}[Duration-blind bookkeeping inflates long commitments]
\label{prop:bookkeeping}
For every $s$, $a_{1:H}$, and $k\in\{1,\dots,H\}$,
\[
\widetilde Q(s,a_{1:H},k)-Q(s,a_{1:H},k)
  \;=\;\bigl(\gamma-\gamma^{k}\bigr)\,\mathbb E\bigl[V(s_k)\bigr]\;\ge\;0 ,
\]
with equality iff $k=1$ or the landing value vanishes. The distortion is strictly increasing in
$k$ and proportional to the value of the landing state.
\end{proposition}
\begin{proof}
Subtract \eqref{eq:smdp-backup} from $\widetilde Q$; all terms cancel except the bootstrap, and
$\gamma\ge\gamma^{k}$ with $V\ge0$ since rewards are nonnegative.
\end{proof}

Under goal-reaching reward the distortion is easiest to see. A $k$-step commitment that reaches
the goal is valued $\gamma$ instead of $\gamma^{k}$, as if the commitment were a single-step
shortcut through the decision graph. Any commitment rule trained or evaluated against
$\widetilde Q$ therefore receives a spurious monotone pressure toward long commitments that has
nothing to do with task outcomes. This failure mode of naive backups for variable execution
horizons was identified in the online setting by ExRL~\citep{exrl}, whose remedy is likewise the
semi-Markov backup. Our critic trains on \eqref{eq:smdp-backup} directly, so commitments of
different lengths compete on equal terms; every result below assumes this bookkeeping.

\subsection{Selection. What a value applied around a frozen policy can and cannot do}
\label{app:selection}

\begin{proposition}[Support ceiling and winner's curse]
\label{prop:selection}
Fix a state $s$ and a frozen sampler $\bar\pi$. Let $a^{(1)},\dots,a^{(N)}\sim\bar\pi(\cdot\mid s)$
be independent proposals, let $\widehat Q=Q+\epsilon$ be the critic with zero-mean,
$\sigma$-sub-Gaussian errors $\epsilon_i$ independent across proposals, and let
$\hat a=\arg\max_i\widehat Q(s,a^{(i)})$ be the selected chunk. Then
\begin{enumerate}
\item[(i)] the realized value obeys
$Q(s,\hat a)\le\sup_{a\in\operatorname{supp}\bar\pi(\cdot\mid s)}Q(s,a)$ almost surely, for every
$N$; and
\item[(ii)] the believed value of the selection overstates it,
\[
\mathbb E\bigl[\widehat Q(s,\hat a)\bigr]-\mathbb E\bigl[Q(s,\hat a)\bigr]
\;\ge\;0,
\qquad
\mathbb E\bigl[\widehat Q(s,\hat a)\bigr]
\;\le\;
\mathbb E\Bigl[\max_i Q(s,a^{(i)})\Bigr]+\sigma\sqrt{2\ln N}.
\]
\end{enumerate}
\end{proposition}
\begin{proof}
(i) is immediate since $\hat a$ lies in the support. For (ii), the first inequality follows from
$\widehat Q(s,\hat a)=\max_i(Q+\epsilon)_i\ge Q(s,\hat a)+\epsilon_{\hat a}$ and
$\mathbb E[\epsilon_{\hat a}\mid\hat a]$ being the error of the argmax, which is nonnegative in
expectation because larger errors are more likely to be selected; formally
$\mathbb E[\max_i(Q_i+\epsilon_i)]\ge\mathbb E[Q_{\hat a}+\epsilon_{\hat a}]$ with
$\mathbb E[\max_i(Q_i+\epsilon_i)]\le\mathbb E[\max_iQ_i]+\mathbb E[\max_i\epsilon_i]$ and the
standard sub-Gaussian maximal inequality $\mathbb E[\max_i\epsilon_i]\le\sigma\sqrt{2\ln N}$.
\end{proof}

Part (i) is the ceiling that best-of-$N$ execution interpolates toward as $N$ grows~\citep{emaq};
part (ii) says that scaling $N$ widens the gap between what the system believes and what it
achieves, since the believed value grows like $\sigma\sqrt{2\ln N}$ while the realized value is
capped by (i). Neither defect is repaired by more compute at deployment.

Selection also cannot move the commitment distribution. The next lemma makes precise why
improvement must be applied to the full chunk.

\begin{lemma}[The full-chunk objective is a tight lower bound on deployed value]
\label{lem:fullchunk}
Let $\kappa^{*}_{\pi}(s)\in\arg\max_{k}Q^{\pi}(s,\mu_\pi(s),k)$ be the selector induced by exact
option values, where $\mu_\pi(s)$ denotes the deployed chunk. Then for every $\pi$,
\[
V^{\pi,\kappa^{*}_{\pi}}\;\ge\;V^{\pi,H},
\]
so the full-chunk objective $J_H(\pi):=\mathbb E_{s}\bigl[Q^{\pi}(s,\mu_\pi(s),H)\bigr]$ is a lower
bound on the deployed value, and the bound is tight exactly when adaptive execution no longer
helps, that is when $V^{\pi,\kappa^{*}_{\pi}}=V^{\pi,H}$.
\end{lemma}
\begin{proof}
$\kappa\equiv H$ is among the candidates the exact selector maximizes over.
\end{proof}

\begin{remark}[Selection-only systems inherit their commitment distribution]
\label{rem:frozen}
If the policy is frozen and only $\kappa$ is tuned, the option values
$Q^{\bar\pi}(\cdot,\cdot,k)$ are fixed, so the selected commitments are a deterministic functional
of the initial policy and never lengthen as competence would allow. Conversely, improving only the
selected short prefixes raises $Q(\cdot,\cdot,k)$ for small $k$ while leaving the full-chunk value
unimproved, so by Lemma~\ref{lem:fullchunk} the deployed lower bound does not move and no
commitment growth occurs. Improvement must therefore attack the full chunk, with the selector
harvesting whatever reactivity remains valuable.
\end{remark}

\subsection{Reactivity. Decomposition, absorption, and the floor}
\label{app:reactivity}

\begin{theorem}[Decomposition of the shortfall]
\label{thm:decomp}
For any chunk policy $\pi$,
\[
\underbrace{V^{\star}_{1}-V^{\pi,H}}_{\text{total shortfall}}
=\underbrace{\bigl(V^{\star}_{1}-V^{\star}_{H}\bigr)}_{\Delta_{\mathrm{alea}}\;\text{(policy independent)}}
+\underbrace{\bigl(V^{\star}_{H}-V^{\pi,H}\bigr)}_{\Delta_{\mathrm{epis}}(\pi)\;\text{(vanishes under improvement)}} .
\]
\end{theorem}
\begin{proof}
Add and subtract $V^{\star}_{H}$. The first summand does not depend on $\pi$; the second is
nonnegative and has infimum $0$ over $\pi$ by definition of $V^{\star}_{H}$.
\end{proof}

\begin{theorem}[Determinism makes reactivity worthless]
\label{thm:determinism}
If $T$ and $r$ are deterministic and the state is fully observed, then
$V^{\star}_{1}=V^{\star}_{H}=V^{\star}_{\mathrm{ada}}$.
\end{theorem}
\begin{proof}
From any $s$ the closed-loop optimum traces a single trajectory. Define $\tilde\pi(s)$ to emit its
first $H$ actions; $\tilde\pi$ is an admissible chunk policy. Executed with full commitment,
determinism delivers exactly the state the optimal trajectory visits at time $H$, where requerying
continues the optimal trajectory by the Markov property. By induction the full-commitment
execution of $\tilde\pi$ reproduces the optimal trajectory, so $V^{\tilde\pi,H}=V^{\star}_{1}$ and
hence $V^{\star}_{H}\ge V^{\star}_{1}$; Lemma~\ref{lem:sandwich} closes the chain.
\end{proof}

\begin{corollary}[Absorption]
\label{cor:absorb}
Under the hypotheses of Theorem~\ref{thm:determinism}, for every $(\pi,\kappa)$,
\[
V^{\pi,\kappa}-V^{\pi,H}\;\le\;V^{\star}_{\mathrm{ada}}-V^{\pi,H}\;=\;\Delta_{\mathrm{epis}}(\pi).
\]
Every gain that adapting the commitment can produce is achievable instead by a better full chunk;
policy improvement absorbs it.
\end{corollary}

\begin{remark}[Reactive information]
\label{rem:reactive-info}
The proof of Theorem~\ref{thm:determinism} uses determinism only through the fact that no
information unpredictable from $s_t$ arrives during $(t,t+H]$. The value of reactivity is
therefore nonzero exactly when such information arrives, through stochastic dynamics or through
partial observability. Vision-based manipulation has both, notably occlusion and unobserved
contact forces, so the floor below is positive in our setting and its size is a measurable
property of each task.
\end{remark}

\begin{theorem}[Floor and gain decomposition]
\label{thm:floor}
Let $\Delta_{\mathrm{react}}:=V^{\star}_{\mathrm{ada}}-V^{\star}_{H}\ge0$. Then
$\Delta_{\mathrm{react}}\le\Delta_{\mathrm{alea}}$, and for every $(\pi,\kappa)$,
\[
V^{\pi,\kappa}-V^{\pi,H}\;\le\;\Delta_{\mathrm{react}}+\Delta_{\mathrm{epis}}(\pi).
\]
Under determinism $\Delta_{\mathrm{react}}=0$ by Theorem~\ref{thm:determinism}, and in general
$\Delta_{\mathrm{alea}}$ is bounded by the reactivity penalty of open-loop
execution~\citep[Cor.~1]{dqc}.
\end{theorem}
\begin{proof}
$\Delta_{\mathrm{react}}\le\Delta_{\mathrm{alea}}$ follows from
$V^{\star}_{\mathrm{ada}}\le V^{\star}_{1}$ (Lemma~\ref{lem:sandwich}). For the gain bound,
$V^{\pi,\kappa}\le V^{\star}_{\mathrm{ada}}=V^{\star}_{H}+\Delta_{\mathrm{react}}$ and
$V^{\star}_{H}-V^{\pi,H}=\Delta_{\mathrm{epis}}(\pi)$.
\end{proof}

The two summands have different fates. $\Delta_{\mathrm{epis}}$ is removed by improving the
policy, and by Corollary~\ref{cor:absorb} this removal also absorbs the corresponding share of the
adaptive gain; $\Delta_{\mathrm{react}}$ can only ever be collected at execution time, by
committing for shorter horizons where reactive information genuinely arrives. A method that only
selects collects at most $\Delta_{\mathrm{react}}$ and forfeits $\Delta_{\mathrm{epis}}$
(Remark~\ref{rem:frozen}); a method that improves at fixed full commitment absorbs
$\Delta_{\mathrm{epis}}$ and forfeits $\Delta_{\mathrm{react}}$. Joint optimization is the unique
combination that takes both.

\begin{corollary}[Curriculum]
\label{cor:curriculum}
As improvement drives $\Delta_{\mathrm{epis}}(\pi)\to0$, the set of states on which a strictly
shorter commitment beats the full chunk shrinks to the states where reactive information arrives,
so the mean selected commitment increases toward a limit set by the environment alone. No
replanning cost is needed to produce this growth; where values tie near convergence, a
lexicographic preference for the longest commitment within an $\epsilon$-optimal set resolves the
indifference with $\epsilon$ its only free parameter.
\end{corollary}

\subsection{Leakage. Closed-loop demonstrations inflate long commitments in a second way}
\label{app:leakage}

Teleoperated demonstrations are closed loop; the recorded action at each step reacts to the
realized outcome of the previous ones. Conditioning a chunk-level regression target on the
executed segment $a_{1:k}$ therefore conditions on outcome information that the demonstrator's
reactions reveal, and the estimated value of executing that segment open loop inherits an optimism
that the open-loop executor cannot realize.

\begin{fact}[\citealp{dqc}]
\label{fact:dqc}
There exist MDPs and closed-loop data distributions on which the chunk-conditioned value estimate
prefers a chunk whose actual open-loop value is worse by $\Omega\!\left(1/(1-\gamma)\right)$, and
under an $\varepsilon_h$-open-loop-consistency assumption the estimation bias admits an upper
bound that is monotonically increasing in the committed length.
\end{fact}

\begin{proposition}[The prefix baseline does not cancel leakage]
\label{prop:leakage}
Write $\widehat Q_k=Q_k+b^{Q}_{k}$ and $\widehat V_k=V_k+b^{V}_{k}$ for the estimated option value
and its state baseline at commitment $k$, and consider any commitment score of the normalized
advantage form $(\widehat Q_k-\widehat V_k)/\gamma^{k}$. Then
\[
\frac{\widehat Q_k-\widehat V_k}{\gamma^{k}}
=\underbrace{\frac{Q_k-V_k}{\gamma^{k}}}_{\text{true normalized advantage}}
+\underbrace{\frac{b^{Q}_{k}-b^{V}_{k}}{\gamma^{k}}}_{\text{selection bias}},
\]
where (i) $b^{Q}_{k}$ contains the closed-loop leakage of Fact~\ref{fact:dqc} and grows with $k$;
(ii) $b^{V}_{k}$ does not contain the corresponding component, because the baseline conditions on
the state alone and not on the executed chunk, so the difference does not vanish; and (iii) the
$\gamma^{-k}$ normalization amplifies the residual exponentially in $k$.
\end{proposition}

The two biases that act on the commitment therefore have different characters and different
cures. The bookkeeping bias of Proposition~\ref{prop:bookkeeping} is a property of the estimator
and is removed exactly by the backup \eqref{eq:smdp-backup}. The leakage bias is a property of the
data-generating process; it survives correct bookkeeping, grows with $k$, and pushes in the
opposite direction to the discount-scale artifact, so correcting only one of the two amplifies the
other. This is why our targets are built conservatively within the support of the
data~\citep{iql, deas} rather than by trusting chunk-conditioned estimates at face value, and why
the commitment must be learned rather than scored.

\subsection{A worked example where the optimal commitment is state dependent}
\label{app:example}

The forces above already imply that no fixed commitment is optimal in general; a two-situation
example makes the separation strict and quantitative. Let the chunk length be $H=2$ and let the
policy resample its chunk imperfectly at every query, producing a correct chunk with probability
$1-\varepsilon$ and an arbitrary one otherwise.

In a \emph{corridor} the dynamics are deterministic and a correct two-step sequence succeeds.
Committing fully succeeds with probability $1-\varepsilon$, since one query suffices; requerying
after one step succeeds with probability $(1-\varepsilon)^{2}$, since each query risks a bad
resample. Commitment wins by $\varepsilon(1-\varepsilon)>0$.

In a \emph{junction} the environment reveals a binary event after the first step and the correct
second action depends on it. Committing fixes the second action before the event and succeeds with
probability $\tfrac12(1-\varepsilon)$; requerying observes the event and succeeds with probability
$(1-\varepsilon)^{2}$, which is larger whenever $\varepsilon<\tfrac12$. Reaction wins.

For every $\varepsilon\in(0,\tfrac12)$ any constant commitment is strictly suboptimal in one of
the two situations, while a state-dependent $\kappa$ takes both, and the margin in the corridor
shrinks as $\varepsilon\to0$, which is the curriculum of Corollary~\ref{cor:curriculum} in
miniature. The junction gain persists, which is the floor of Theorem~\ref{thm:floor}.

\subsection{From theory to measurement}
\label{app:predictions}

Each result above yields a falsifiable prediction that our experiments test; none of them requires
the full method to run.

\begin{center}
\begin{tabular}{@{}llll@{}}
\toprule
Result & Prediction & Measurement & Rejected if \\
\midrule
Prop.~\ref{prop:bookkeeping} & duration-blind scoring shifts & same critic, two scorers, & $k$ histograms \\
 & commitments long and loses & $k$ histogram and success & indistinguishable \\
\addlinespace
Prop.~\ref{prop:selection} & best-of-$N$ plateaus while its & $N$ sweep; believed minus & realized value keeps \\
 & believed value keeps rising & realized value of selection & pace with believed \\
\addlinespace
Fact~\ref{fact:dqc}, Prop.~\ref{prop:leakage} & chunk-value optimism grows & open-loop replay of demo & gap flat in $k$ \\
 & with committed length & prefixes vs.\ critic estimate & \\
\addlinespace
Thm.~\ref{thm:floor} & best fixed $k$ varies by task; & fixed-$k$ sweep plus oracle & oracle matches best \\
 & state-dependent $\kappa$ beats it & per-state selection & fixed $k$ everywhere \\
\addlinespace
Cor.~\ref{cor:curriculum} & mean commitment grows only & improvement on/off arms, & growth without \\
 & under policy improvement & mean selected $k$ over training & improvement \\
\bottomrule
\end{tabular}
\end{center}

The first three rows probe the estimator and require only the critic and the demonstrations; the
last two probe execution and require rollouts of the base policy. Section~\ref{sec:experiments}
reports all five, with protocols registered before the runs.
